\section{The actual base theta map: convergence, descent, and the remaining extension}
\label{sec:base-theta-v032}

\status{BTH1--BTH15 give complete proofs of the stated implications on
the original base tuple.  They produce its convergent lowest-type
integral, reconstruct the smooth-vector uniqueness and finite
orthogonal-kernel arguments, and reduce the original base theta square
to an explicit common-object producer.  They do not yet construct the
original smooth cuspidal extension, its geometric reading, or its
simultaneous rational tests.  Those producer obligations remain C.
No new Lean replay is claimed.}

The base-rank consumer in SLR8 needs \(\mathrm{RowsHold0}_3\) and a
theta square at \(t_0=A_3.A.\mu_0\).  Its need for that square is
different from the separate \(\mathrm{RankM}\) branch.  Everything
below concerns the former consumer and one fixed original bundle.
There is no change to its coefficient field, Hilbert metric, adjoint,
finite group action, or chosen theta map.

Here is the notation used throughout.  Put
\[
 K=D.\mathrm{what}_2.\mathrm{LL},\qquad
 \iota=D.\mathrm{what}_2.\mathrm{embLL},\qquad
 \sigma=A_3.A.v_0.\mathrm{embedding}.
 \tag{BTH.1}
\]
Let \(R,F,Y,Z\) be the original real datum, face, \texttt{Read},
and \texttt{Read\allowbreak L}; their identification with the supplied bundle is
an explicit premise.  For \(i<m\), put
\[
 \begin{aligned}
 H_i&=(R.HV(i,t_0)).H,&
 W_i&=(R.HW(i,\mathrm{sg}_0)).H,\\
 L_i&=R.\mathrm{extV}(i,t_0),&
 U_i&=R.\mathrm{thF}(i,t_0),&
 T_i&=R.\mathrm{thG}(i,t_0).
 \end{aligned}                                                  \tag{BTH.2}
\]
The realization \(e_i:W_i\to\mathrm{MMap}(\mathcal D_i)\) is injective
and has image \(\mathrm{Wall}_i\);
\(\mathcal D_i=R.MW(i,\mathrm{sg}_0)\) is the prescribed tensor of
discrete-series models.  These are
fields of \texttt{Is\allowbreak Models}.  The space \(\mathrm{MMap}(\mathcal D_i)\)
consists of multilinear maps from the smooth-vector factors to
automorphic functions.  The subspace \(\mathrm{Wall}_i\) is the
span of maps with the source's automorphy, cuspidality, level,
equivariance, and Sobolev-continuity properties.

The full level coordinate space will be denoted \(K^{d_i}\).
Its map to rational forms is
\[
 f_i(x)=\mathrm{liftForm}(i,t_0,x).
 \tag{BTH.3}
\]
The original good-cut space is \(K^{s_i}\), and its inclusion and
normalized projection are
\(E_i=lZ0(i,t_0)\) and \(\Pi_i=lPrj(i,t_0)\).
The original target-coordinate map is \(J_i=Z.jW(i)\), and its
Hilbert reading is \(q_i=Y.mj.eW(i)\).
In particular \(q_iJ_i=\mathrm{goodForm}_i\).  SLR2 proves the
normalization
\[
 \Pi_iE_i=1,\qquad
 x-E_i\Pi_i x\perp_{lH_i}\operatorname{im}E_i.                    \tag{BTH.4}
\]
The tuple identification is literal:
\(Y.\mathrm{tup}(t_0)=D.\mathrm{what}_2.\mu_0\), which is
always \(\mathrm{Used0}\).  No assertion about a nonbase tuple
is needed.

% BEGIN BTH1
\paragraph{BTH1: smooth vectors of the lowest compact type are nonzero and cyclic in the required topology.}
\label{claim:BTH1}
Let \(\pi\) be one of the source's nonzero irreducible unitary
representations of a real unitary Lie group \(G_\infty\), with
maximal compact subgroup \(K_\infty\).  Write \(H^\infty\) for its
smooth vectors with seminorms
\[
 p_{\mathcal X}(v)=\|d\pi(X_1)\cdots d\pi(X_s)v\|,
 \quad X_j\in\mathfrak g,\quad s\ge0.                            \tag{BTH.5}
\]
Its nonzero lowest compact isotypic space contains a nonzero
smooth vector.  Moreover, if \(0\ne u\in H^\infty\), every
continuous linear functional for the topology (BTH.5) that
vanishes on all \(\pi(g)u\) is zero.

\emph{Proof.}
Choose a nonnegative smooth bump \(\phi\) supported sufficiently
near \(1\), and normalize its Haar integral to one.  Averaging
\(\phi(k g k^{-1})\) over \(K_\infty\) makes it conjugation
invariant under \(K_\infty\), with arbitrarily small support.
Conjugation preserves Haar measure here: the modular function
has compact image on \(K_\infty\), so that image is \(\{1\}\).
The smoothing operator
\[
       S_\phi v=\int_{G_\infty}\phi(g)\pi(g)v\,dg
 \tag{BTH.6}
\]
commutes with \(K_\infty\).  Differentiating the compactly
supported scalar kernel after a change of variable gives
\(S_\phi H\subset H^\infty\).  More precisely, for each
derivative word \(\mathcal X\) there is a compactly supported
smooth derivative \(D_{\mathcal X}\phi\), with the appropriate
signs, such that
\[
 p_{\mathcal X}(S_\phi v)
       \le \|D_{\mathcal X}\phi\|_{L^1}\|v\|.                    \tag{BTH.7}
\]
This follows first for one derivative by differentiating
\(\phi(\exp(-tX)g)\), and then by induction.  Compact support
justifies every differentiation under the integral.
For a nonzero vector \(v\) in the lowest compact isotypic space,
strong continuity permits the support to be chosen so that
\(\|\pi(g)v-v\|<\|v\|/2\) there.  Hence
\(\|S_\phi v-v\|<\|v\|/2\); the smoothed vector is nonzero
and remains in the same isotypic space.

Now suppose \(\ell\) has the stated continuity and vanishing
properties.  Its continuity gives finitely many words and a
constant \(C\) with
\[
       |\ell(v)|\le C\sum_{\mathcal X}p_{\mathcal X}(v).
 \tag{BTH.8}
\]
By (BTH.7), \(\ell S_\phi\) extends to a bounded functional on
the Hilbert space \(H\).  For every \(g\),
\[
  \ell S_\phi(\pi(g)u)
      =\int\phi(h)\ell(\pi(hg)u)\,dh=0.                          \tag{BTH.9}
\]
The interchange is valid in the smooth topology: for a fixed
smooth \(u\), all derivative seminorms of the orbit map are
continuous and bounded on this compact support.
The closed Hilbert span of \(\{\pi(g)u:g\in G_\infty\}\) is a
nonzero invariant subspace, hence all of \(H\) by irreducibility.
Thus \(\ell S_\phi=0\).

Take bumps forming an approximate identity.  For every
\(v\in H^\infty\), \(S_\phi v\to v\) in every seminorm (BTH.5).
To see this directly, move a derivative across \(\pi(g)\) using
\(d\pi(X)\pi(g)=\pi(g)d\pi(\operatorname{Ad}(g^{-1})X)\);
iteration gives a finite sum whose coefficients and vectors
converge as \(g\to1\).  Integrating over shrinking supports gives
the convergence in each fixed seminorm.  Consequently
\(\ell(v)=\lim\ell(S_\phi v)=0\).

The source's \texttt{DSModel.ds} supplies the irreducibility
and unitarity used here, and its lowest compact type is nonzero
before smoothing.  Its frames conjugate the standard unitary
group to the original Gram-matrix group, preserving the argument.
This proves the smooth-topology statement; Hilbert density alone
would not have justified applying a Sobolev-continuous functional.
The corresponding source declarations are
\texttt{exists\_\allowbreak smooth\allowbreak Vec\_\allowbreak low\allowbreak KSub}, \texttt{low\allowbreak T\_\allowbreak ne\_\allowbreak bot},
\texttt{eq\_\allowbreak zero\_\allowbreak of\_\allowbreak translates}, and \texttt{low\allowbreak Generates}.
% END BTH1

% BEGIN BTH2
\paragraph{BTH2: a multiplicity map is determined by its lowest-type values.}
\label{claim:BTH2}
For the original \(\mathcal D_i\), suppose \(A,B\in\mathrm{Wall}_i\)
have equal values on all pure tensors of lowest-type vectors.
Then \(A=B\).

\emph{Proof.}
Put \(C=A-B\).  A scalar evaluation of \(C\), with the other
archimedean variables fixed, is Sobolev-continuous in the
remaining variable.  For a generator of \(\mathrm{Wall}_i\)
this is part of its defining continuity condition.  It remains
true under finite sums and scalar multiples, which is exactly
how the span defining \(\mathrm{Wall}_i\) is formed.

Order the finite set of archimedean factors.  Initially \(C\)
vanishes when every factor is of lowest type.  Suppose it
vanishes with the first \(s\) factors arbitrary and the rest
of lowest type.  Fix all but the next factor, and fix the
adelic evaluation point.  Equivariance in that next factor
shows that the resulting functional vanishes on every
translate of a nonzero smooth lowest-type vector: the translate
only changes the adelic evaluation point in a function that
already vanishes identically.  BTH1 makes this functional
zero on the entire smooth factor.  Induction removes all the
lowest-type restrictions.  Multilinearity and the spanning
property of pure tensors give \(C=0\).

The same proof applies to the matched finite target lifts used
below.  Those lifts normalize the prescribed principal flag
level and preserve its character, by the original level-lift
conditions.  Their right translations therefore preserve this
\(\mathrm{Wall}_i\), including its automorphic and Sobolev
conditions.  Thus an identity proved by the theta
integral on lowest types extends to all smooth-vector
arguments.  This is the exact content needed from
\texttt{eq\_\allowbreak zero\_\allowbreak of\_\allowbreak low\allowbreak T} and
\texttt{theta\allowbreak Lift\_\allowbreak intertwines\_\allowbreak lifts\_\allowbreak smooth\allowbreak Vectors}.
It is a uniqueness theorem, not an existence theorem for the
theta extension.
% END BTH2

% BEGIN BTH3
\paragraph{BTH3: the original base lowest-type theta integral converges.}
\label{claim:BTH3}
Assume the original model fields and weak theta fields
\(\omega P,vP,pt\_theta\).  For every \(y\in H_i\), every
lowest-type tensor \(a\), and every target adelic point \(g\),
the original integral
\[
 {\cal I}_i(y,a)(g)=
 \int_{\mathcal F_{V,i}}
  \theta KO(\omega_i)(h,g)
       \bigl(pt_i(io_i(a,(L_i y)(h)))\bigr)\,d\nu_i(h)            \tag{BTH.10}
\]
is absolutely convergent.  It is complex linear in \(y\)
and in \(a\).  Neither an \texttt{Is\allowbreak Io\allowbreak W} law nor an existing
cuspidal extension is needed for this conclusion.

\emph{Proof.}
The source space \(V_c(D,i)\) is positive definite at every
archimedean place: \texttt{models.vpos}, the original
\texttt{V\_\allowbreak sig}, and a change of frame give positive definite
Gram matrices.  A real unitary group of such a matrix is compact.
Indeed write its Gram matrix as \(B^*B\).  Conjugation by \(B\)
identifies the group with \(U(n)\), whose entries have absolute
value at most one and whose defining equations are closed.

For the finite part use the double-coset representatives
stored in \((R.B.L(i,t_0)).pres\); no new class-number
finiteness theorem is required.  The finite level is a closed
subgroup of a conjugate of \(GL_n(\widehat{\mathcal O}_E)\)
cut out by finitely many congruences.  It is compact: each
integral local matrix group is compact, the product is compact,
and conjugation and the closed unitary equations preserve this
property.  Its inverse image in the full adelic group is compact
because its real factor is compact as well.  Lifting the finitely
many stored representatives and taking their translates gives
a compact set \(C\) with
\[
       G_V(\mathbb A)=G_V(F)\,C.                                \tag{BTH.11}
\]
The finite projection used here is the source's range
projection, so each stored representative has an adelic lift.

The prescribed measure \(\nu_i\) is Haar and its prescribed
\(\mathcal F_{V,i}\) is a measurable fundamental domain.  It
has finite measure.  For completeness, the countable rational
group and (BTH.11) give
\(\sum_{\gamma\in G_V(F)}1_C(\gamma h)\ge1\).
Integrate over that fundamental domain and unfold its
translates.  Countable additivity for the nonnegative sum and
left invariance give
\[
 \nu_i(\mathcal F_{V,i})
 \le\int_{\mathcal F_{V,i}}\sum_\gamma1_C(\gamma h)\,d\nu_i(h)
 =\nu_i(C)<\infty.                                               \tag{BTH.12}
\]
The fundamental-domain equality is understood almost everywhere,
as in its definition.

The function \(F_y=L_i y\) is continuous and rational-left
invariant.  The first assertion follows from the continuous
coefficient representation and a sufficiently small principal
open subgroup on which the finite-level character is one.
The second follows by cancellation in
\[
     (c_V.\mathrm{ext}\,f)(h)
       =\rho_V(h^{-1})f(\mathrm{fin}(h)):
\]
the rational automorphy of \(f=(R.HV(i,t_0)).e(y)\) cancels
\(\rho_V(\gamma)\).

Choose a basis \(b_\alpha\) of the finite-dimensional
coefficient space of \(c_V\).  The integrand in (BTH.10) is
the finite sum
\[
 \sum_\alpha b_\alpha^*(F_y(h))\,
    \omega_i(h,g)\bigl(pt_i(io_i(a,b_\alpha))\bigr)(0).            \tag{BTH.13}
\]
Its first factors are continuous, and its second factors
are continuous by the \texttt{om\_\allowbreak cont} field of the actual
\texttt{Is\allowbreak Weil\allowbreak Datum\allowbreak O}.  The \texttt{pt\_\allowbreak theta} field checks
the exact theta-function hypothesis for each vector appearing
in (BTH.13).  Rational-left invariance follows from
\texttt{om\_\allowbreak aut}, \texttt{om\_\allowbreak pres}, and the invariance of
\(F_y\).  It is therefore bounded on all of the quotient by
its bound on \(C\).  Equations (BTH.11)--(BTH.12) prove absolute
integrability.  Bilinearity of \(io_i\), linearity of the
Weil action and \(pt_i\), and additivity of integrals of
integrable functions prove the two linearity assertions.

Thus the \texttt{Is\allowbreak Theta\allowbreak P} integrability guard is discharged
on the actual base data.  The substantially stronger assertion
that (BTH.10) extends to a smooth cuspidal multiplicity map is
not a consequence of compactness of the source quotient.
% END BTH3

% BEGIN BTH4
\paragraph{BTH4: finite extension data activate the original weak choices and construct the good-projected map.}
\label{claim:BTH4}
Fix a finite basis \(h_1,\ldots,h_n\) of the original \(H_i\).
Assume that for each \(h_\alpha\) there is
\(A_\alpha\in\mathrm{Wall}_i\) whose lowest-type restriction
is the integral (BTH.10).  Then the source's actual stored
\(U_i=\mathrm{thF}(i,t_0)\) satisfies \(\mathrm{PThF}(i,t_0)\).
The actual stored \(T_i=\mathrm{thG}(i,t_0)\) then satisfies
\[
       e_i T_i=\mathrm{postOn}(eWf_i)\,U_i.                      \tag{BTH.14}
\]
No independent \(\mathrm{PThG}\)-existence premise is needed.

\emph{Proof.}
Define \(U(\sum c_\alpha h_\alpha)=\sum c_\alpha A_\alpha\).
It is linear and lies in \(\mathrm{Wall}_i\).  BTH3 and the
basis expansion identify its lowest-type restriction with
\({\cal I}_i\) for every \(y\).  On \(\operatorname{im}L_i\),
use \(L_i^{-1}\), available from \(Y.\mathrm{hinj}\), to define
the low-type theta map.  Formula (BTH.10) and its convergence
give exactly the two fields of \texttt{Is\allowbreak Theta\allowbreak P}.
This proves \(\mathrm{PThF}(i,t_0,U)\).
BTH2 also proves uniqueness of such \(U\).

The weak-choice rule is
\[
 \mathrm{OrZero}(P,z)
 =((\exists w,\ P(w))\Rightarrow P(z))
   \ \wedge\ ((\nexists w,\ P(w))\Rightarrow z=0).                \tag{BTH.15}
\]
The first implication in \texttt{tw\_\allowbreak th\allowbreak F} therefore gives
\(\mathrm{PThF}(i,t_0,U_i)\) for the stored \(U_i\), not
for a substituted map.

Write the prescribed central projector as
\(eWf_i=\sum_{g\in G}c_g\,\mathrm{actL}(g)\).
The model identity for \(gW\) gives
\[
 e_i\!\left(\sum_g c_g\,gW_i(g)w\right)
       =\mathrm{postOn}(eWf_i)(e_iw).                           \tag{BTH.16}
\]
Since \(\operatorname{im}e_i=\mathrm{Wall}_i\), the map on
the right of (BTH.14) takes its values in
\(\operatorname{im}e_i\).  Pulling it back through the linear
isomorphism \(e_i:W_i\to\operatorname{im}e_i\) constructs a
linear map \(T\) satisfying (BTH.14).  Injectivity of \(e_i\)
proves its uniqueness.  The first implication in
\texttt{tw\_\allowbreak th\allowbreak G} now proves (BTH.14) for the stored \(T_i\).

This is the pointwise-at-\(t_0\) form of
\texttt{exists\_\allowbreak linear\_\allowbreak lift} and \texttt{exists\_\allowbreak factor\_\allowbreak hmod}.
The global eight-field \texttt{Ex\allowbreak At} package, including its
chain lifts and other tuples, has not been used.
% END BTH4

% BEGIN BTH5
\paragraph{BTH5: the actual finite character gives theta covariance and the inverse target level.}
\label{claim:BTH5}
Assume the original weak Weil, polynomial, group, and level
fields and the base extension from BTH4.  For matching source
and target finite-level lifts \(\kappa,\kappa'\), put
\(c=\mathrm{levChi}(\mathrm{flagL}(\mathrm{datV}))(\kappa)\).
Then
\[
 \omega_i(\kappa,\kappa')vP_i=c^{-1}vP_i.                        \tag{BTH.17}
\]
If
\[
        c\,L_i y'(h)=L_i y(h\kappa),
\]
the full multiplicity maps satisfy
\[
        U_i y'(v)(g)=U_i y(v)(g\kappa')                         \tag{BTH.18}
\]
for every smooth-vector tuple \(v\).  In particular a source
level character is inverted on the target level, with the
source's exact fixed embedding.

\emph{Proof.}
We record the finite normalization because a scalar ambiguity
would change the asserted map.  At a prime in \(P_g\), write
\(a,d\) for the determinant characters of the source's upper
and lower blocks and \(a'\) for the target character.  The
matched upper blocks give \(a'=a\).  At a prime in \(Q\), write
\(x_a,x_d,x_b\) for the corresponding prescribed \(\xi\)-values.
The source's vacuum and half-datum transformations are
\[
 \begin{aligned}
 m_V&=\prod_{P_g}(ad)^{-1}\prod_Q x_ax_d,&
 m_W&=\prod_{P_g}(a')^{-1}\prod_Q x_b,\\
 d_{\rm half}&=\prod_Q x_a^{-1}x_b,&
 c&=\prod_{P_g}d^{-1}\prod_Q x_d .
 \end{aligned}                                                  \tag{BTH.19}
\]
Thus \(m_V^{-1}m_W=d_{\rm half}c^{-1}\).
The half-support is preserved: its reduced adjoint frame
changes by \(A\mapsto GAH^{-1}\); the zero lower block
remains zero, the identity upper block at \(P_g\) remains
the identity, and the invertible upper block at \(Q\)
has inverse-determinant-character multiplier
\(x_a^{-1}x_b\).  Reindexing the finite Heisenberg sum therefore
gives
\(d_{\rm half}\omega_i(\kappa,\kappa')vP_i
   =(m_V^{-1}m_W)vP_i\).
All factors are units, so cancellation proves (BTH.17).
The vacuum multipliers are precisely the actual-action
stabilizer character proved in BTH15 below, with the recorded
determinants.  BTH13--BTH14 justify the identification with a
canonical action wherever that comparison is used.  Their
frame hypotheses are the self-dual
\(\beta V,\beta W\) fields in the original level data.
Their unramified-away hypothesis is the original
\(\mu\)-character condition away from \(P_g\cup Q\);
the second Weil character is \(\mu\eta^{-1}\).
The required extension character is the original \(\eta\):
the two restriction laws and
\(\operatorname{rank}V=r+t,\operatorname{rank}W=r\),
with \(t\) odd, give its required odd-parity restriction.
Explicitly, for every quadratic \(\eta_F\) appearing in
\texttt{Res\allowbreak Eta}, the two restrictions give
\(\mu|_F=\eta_F^{r+t}\) and
\((\mu\eta^{-1})|_F=\eta_F^r\).  Since
\(\eta_F^2=1\), division gives
\(\eta|_F=\eta_F^{2r+t}=\eta_F\).
This holds for every such \(\eta_F\), retaining the universal
restriction in \texttt{Is\allowbreak Eta\allowbreak Tilde}.
The level record's \texttt{Chr\allowbreak Eq}/\texttt{Is\allowbreak Eta\allowbreak T} gives the
prescribed infinity exponent and triviality on each finite
local unit group.  Products with finite support are dense in
the compact product of these groups, so continuity gives
triviality on every finite integral unit idele.
These are the checks carried out in
\texttt{om\_\allowbreak lifts} before its finite determinant cancellation.
No arbitrary Weil scalar is introduced.

The finite actions commute with the polynomial creation
operators at infinity.  This differentiation does not
presuppose continuity of an ambient endomorphism on all
functions: the actual Haar-norm identity makes it an isometry
on theta \(L^2\); on a fixed smooth finite level the real
Heisenberg difference quotients converge in \(L^2\), and
covariance carries those quotients through the isometry.
The output is again a smooth theta function, so the resulting
\(L^2\) derivative is its classical derivative.  Iteration is
exactly the difference-quotient argument in WA3.
Applying those operators to
(BTH.17), one at a time, gives the same scalar \(c^{-1}\)
on every \(pt_i(P)\).  Equivalently this is the pure-tensor
finite-action calculation in \texttt{om\_\allowbreak pure\allowbreak Tensor\allowbreak O\_\allowbreak of\_\allowbreak finite}.
The homomorphism law for \(\omega_i\) then gives the kernel
identity
\[
 \theta(g\kappa',h\kappa,pt_i(P))
       =c^{-1}\theta(g,h,pt_i(P)).                              \tag{BTH.20}
\]
Every level lift lies in a compact subgroup of the source
adelic group.  Its modular character is therefore one, so
right translation by \(\kappa\) preserves Haar measure.
The integral of an integrable rational-left-invariant
function over a fundamental domain is unchanged by that
translation: translate the domain, use the fundamental
domain partition, and cancel the left rational translates.
BTH3 checks the integrability for the present functions.
Substitution in (BTH.10), followed by (BTH.20), proves
(BTH.18) on lowest-type tensors.  BTH2 proves it on all
smooth tensors.

To check the inverse character explicitly, let \(p\) be a
target level element.  Its matched source lift follows
from the already supplied \texttt{Bridge.Act\allowbreak Tie.red\_\allowbreak surj}:
choose \(q\) above \(eG^{-1}(g_0)\), lift its finite
component with identity infinite component, and use
\texttt{grp\_\allowbreak red\_\allowbreak L} to identify its reduction with the
literal \(rTup\) of \(g_0\).  Thus no separate
parahoric-surjectivity producer is needed here.
The original level relation says that the source character
on this lift is \(c\,\psi(p)^{-1}\).  Hence one can take
\(y'=\psi(p)^{-1}y\) in (BTH.18), which gives
\(U_i y(v)(gp)=\psi(p)^{-1}U_i y(v)(g)\).
The source declaration
\texttt{lift\allowbreak Level\_\allowbreak fixed\allowbreak Embedding\_\allowbreak eq\_\allowbreak conj\_\allowbreak row\allowbreak Level}
identifies this with \(R.\mathrm{levL}(F,\mathrm{false},
\mathrm{sg}_0)\).  This verifies the full finite-level
membership needed below, not just principal-level membership.
% END BTH5

% BEGIN BTH6
\paragraph{BTH6: the stored target finite action is unitary and the base theta map intertwines it.}
\label{claim:BTH6}
Under the same assumptions, the target action \(gW_i\)
preserves the original Hilbert inner product, and
\[
 T_i jV_i(r_i(g)f)=gW_i(g)\,T_i jV_i(f)
       \qquad(g\in G,\ f\text{ an original rational form}).      \tag{BTH.21}
\]
The central projector acts as the identity on
\(\operatorname{im}T_i\).  Consequently the lift portions of
the source's \texttt{PGUnit}, \texttt{PGEquiv}, and
\texttt{e\allowbreak H\_\allowbreak laws} introduce no additional base producer.

\emph{Proof.}
The Hilbert realization identifies the inner product with
the prescribed multiplicity inner product \(ipAt_i\).
For a finite group element with a target adelic lift, its
action is a right translate multiplied by a level character
of absolute value one.  Right translation preserves Haar
measure because the lift belongs to a compact subgroup.
The defining integral for each multiplicity pairing is
unchanged by the translation and the two unit-modulus
factors cancel.  This also preserves the assertion that a
scalar is the multiplicity pairing.  Its uniqueness is
explicit: choose a nonzero smooth vector in each discrete-series
factor using BTH1, take the two tensor arguments equal, and
divide by the nonzero product of their squared norms.
Translation by the inverse gives the converse existence
implication.  Thus the translation preserves the source's
chosen \(ipAt_i\), including its conditional-choice branch.
Norms, and then inner
products by complex polarization, are unchanged.

If an element has no target lift, the source definition
sets its function action to zero.  The model identity
\(e_i gW_i(g)=\mathrm{postOn}(\mathrm{actL}(g))e_i\)
and injectivity of \(e_i\) would then make \(gW_i(g)\)
zero.  It is an invertible group action, so \(W_i=0\);
all the asserted identities in this case follow.
This deals with the no-lift branch rather than assuming
surjectivity of target lifts.

For a rational form \(f\), the original
\texttt{act\_\allowbreak tie.red\_\allowbreak surj} and \texttt{grp\_\allowbreak red\_\allowbreak L},
by the construction just given in BTH5, supply a source lift \(\kappa\)
with
\[
 (HV.e)(jV_i(r_i(g)f))(h)
 =c^{-1}(HV.e)(jV_i f)(h\,\mathrm{fin}(\kappa)).
 \tag{BTH.22}
\]
The coefficient extension converts (BTH.22) into the
hypothesis of BTH5.  Choose a target lift if one exists.
Equation (BTH.18) and the trivial principal-level target
character identify the resulting function with
\(\mathrm{actL}(g)\,U_i(jV_i f)\).

The good projector commutes with these finite actions.
Indeed its group-algebra element is the original central
character idempotent
\[
C=\sum_{\chi\in X}\frac{\chi(1)}{|G|}
                 \sum_{g\in G}\chi(g^{-1})g ;
 \tag{BTH.23}
\]
conjugation invariance of every \(\chi\in X\) proves centrality by
reindexing, and the irreducible character orthogonality
used in the source's \texttt{Ok} makes these distinct character
idempotents mutually orthogonal and gives \(C^2=C\).
The identification of its coefficients with \(eWf_i\)
is \texttt{grp.record}; the embeddings, including their
complex conjugation, are the ones in (BTH.1).
Apply the good projector to (BTH.18), use (BTH.14), and
then use injectivity of \(e_i\).  This proves (BTH.21).

Finally (BTH.16) writes \(T_i y=E_i^{\rm cen}w\) for
some \(w\), where \(E_i^{\rm cen}=\sum c_g gW_i(g)\).
Since \((E_i^{\rm cen})^2=E_i^{\rm cen}\), this projector
fixes \(T_i y\).  The same statements on the twisted
spaces follow by applying the exact \(\mathrm{Tw}\) functor.
The rational source action is unitary for its original
coordinate Hermitian form by the supplied \texttt{Ok}
block data.  Thus both unitary actions and their actual
intertwiner required in BTH9 are available.
% END BTH6

% BEGIN BTH7
\paragraph{BTH7: the exact original Hilbert reading, with an explicit witness.}
\label{claim:BTH7}
Assume that \(Y,Z\) are the actual readings on the original
bundle, and the base extension hypotheses of BTH4 hold.
For every \(x\in K^{d_i}\),
\[
 q_i\!\left(
 D.\theta P\bigl(\mu_0,i+1,D.jMt(\mu_0,i+1)(Y.eMt(i,t_0)x)\bigr)
       \right)=T_i jV_i(f_i(x)).                               \tag{BTH.24}
\]

\emph{Proof.}
Let \(x_M=D.jMt(\mu_0,i+1)(Y.eMt(i,t_0)x)\) and
\(z=jV_i(f_i(x))\).  The \texttt{mt\allowbreak Read} field says exactly
\[
       (Y.IL.jMt(\mu_0,i,x_M)).1=L_i z.                         \tag{BTH.25}
\]
By the range and level condition in the codomain of
\(IL.jMt\), this function is at the prescribed source
base level.  BTH5 therefore puts \(U_i z\) in the
corresponding full target multiplicity space.
The central projector preserves that target level:
its action commutes with the level action and preserves
the prescribed character.  Its idempotence puts the
projected map in the original good target space \(W_c\).

Use the explicit tensor-linear map
\[
                  w=\operatorname{lift}(U_i z).                \tag{BTH.26}
\]
The definition of \texttt{to\allowbreak Lift\allowbreak Real.th\allowbreak C} first inverts
\(L_i\) on its image and then restricts this tensor map to
lowest types.  Equation (BTH.25), \(L_i\)'s injectivity,
and \(Y.\mathrm{tupT\_tup}\) identify that inverse with
this same \(z\).  Consequently \(w\) has precisely the
required lowest-type restriction.  Restricting the
convergent integral from \(\operatorname{im}L_i\) to its
base-level subspace proves \texttt{Is\allowbreak Theta\allowbreak P} for that
restriction: equality is checked on pure tensors and
integrability is inherited from BTH3.

The universally quantified \texttt{Y.tied.lift.theta}
may now be applied to this particular \(w\).  It gives
\[
      IL.jMj(D.\theta P(\mu_0,i+1,x_M))
             =eWf_i\circ w.                                    \tag{BTH.27}
\]
The \texttt{mj\allowbreak Read} field changes the left side into
\(\operatorname{lift}(e_i q_i(D.\theta P(\cdots)))\).
On the right we have
\[
 \operatorname{lift}(\mathrm{postOn}(eWf_i)U_i z)
     =eWf_i\circ\operatorname{lift}(U_i z).
\]
By (BTH.14) this is
\(\operatorname{lift}(e_i T_i z)\).
Injectivity of the tensor-linear equivalence and of \(e_i\)
proves (BTH.24).

This direct choice avoids obtaining an unspecified witness
from \texttt{Theta\allowbreak Fire\allowbreak L} and then identifying it afterwards.
BTH2 is still used in BTH5 to establish its full-level
membership.  The geometric \texttt{Lift\allowbreak Tie.theta} premise
is still part of the actual reading; this argument does
not construct that premise from an arbitrary bundle.
% END BTH7

% BEGIN BTH8
\paragraph{BTH8: one base rational test family supplies the rational adjoint image.}
\label{claim:BTH8}
Use the source's twisted spaces
\[
 V_i=\mathrm{Tw}_{\epsilon}(H_i),\qquad
 \mathcal W_i=\mathrm{Tw}_{\epsilon}(W_i),\qquad
 \Theta_i=(R.\mathrm{toCx3}(F)).\theta'(i,t_0),                  \tag{BTH.28}
\]
where \(\epsilon=\mathrm{twAt}(\iota,v_0)\).  The phase is
one by \texttt{Pins.phs}, so that
\(\mathrm{Tw.out}\,\Theta_i(\mathrm{Tw.mk}\,y)=T_i y\).
The original rational-form realization
\(j_i:F_i\to_\sigma V_i\) spans \(V_i\) and is isometric
for its original \(K\)-Hermitian form.

Suppose there are finitely many complex-linear functionals
\(\lambda_1,\ldots,\lambda_N\) on \(\mathcal W_i\) such that
\[
 \begin{split}
  &\bigl(\lambda_a(w)=0\text{ for all }a\bigr)\Longrightarrow w=0,\\
  &\lambda_a(\Theta_i j_i f)\in\sigma(K)
                   \quad(f\in F_i),\\
  &\lambda_a(w_{0,i}e_k)\in\sigma(K)
                   \quad(k\text{ an original target basis index}).
 \end{split}                                                    \tag{BTH.29}
\]
Here \(w_{0,i}=(R.\mathrm{toCx3}(F)).w0C(\mathrm{false},i)\).
Then, writing
\(Z_i=\operatorname{im}\Theta_i^*\) and
\(Z_{K,i}=j_i^{-1}(Z_i)\),
\[
                   \operatorname{span}_{\mathbb C}j_i(Z_{K,i})
                         =Z_i.                                 \tag{BTH.30}
\]

\emph{Proof.}
Choose a \(K\)-basis of \(F_i\).  Its Gram matrix for the
original nondegenerate \(K\)-Hermitian form is invertible
and remains invertible under \(\sigma\).  By the isometric
realization, this is the Gram matrix of its image vectors,
so they are complex linearly independent.  The full-span
property makes them a complex basis of \(V_i\).
Applying all the tests to its
theta images gives a matrix with entries in \(\sigma(K)\).
Because the tests separate points, its kernel is exactly
\(\ker\Theta_i\) in these coordinates.  Gaussian elimination
over \(K\) gives a \(K\)-basis of that kernel, and its
complex span is the whole complex kernel: the nonzero
minors determining rank remain nonzero under \(\sigma\).

The matrix of the source Hermitian form also has entries
in \(K\), and \(\sigma(\bar a)=\overline{\sigma(a)}\).
Taking the orthogonal complement of the rational kernel
is therefore another system of linear equations over
\(K\), with the same rank after complexification.
Its complex span is
\((\ker\Theta_i)^\perp=\operatorname{im}\Theta_i^*\).
This is (BTH.30), with the original adjoint and metric.

Only the first two lines of (BTH.29) were used here;
the third will give rational target coordinates in
BTH11.  Thus one base-only \texttt{Lift\allowbreak Test\allowbreak B} family
does both jobs.  No tests for another tuple and no
independent \texttt{Lift\allowbreak Test} family are needed.
Existence of the simultaneous separating tests with
these rationality properties remains a substantive
producer; finite-dimensionality alone does not supply
them.
% END BTH8

% BEGIN BTH9
\paragraph{BTH9: the finite unitary idempotent kernel argument.}
\label{claim:BTH9}
Let \(F_K\) be a rational model of a finite-dimensional
complex Hilbert space \(V\), with full realization \(j\).
Let finite-group representations \(\rho,\tau\) on \(V,W\)
be unitary, let \(\Theta:V\to W\) intertwine them, and
assume
\(\operatorname{span}_{\mathbb C}j(Z_K)
  =Z=\operatorname{im}\Theta^*\).
Let \(e,a\in K[G]\), assume \(\rho(e)^2=\rho(e)\), and
assume \(\tau(a)\) fixes \(\operatorname{im}\Theta\).
For \(f\in F_K\), suppose
\[
 ef=af,\qquad
       \langle jz,jf\rangle=0
           \quad(z\in Z_K\cap\operatorname{im}e).
 \tag{BTH.31}
\]
Then \(\Theta(jf)=0\).

\emph{Proof.}
The kernel of \(\Theta\) is group invariant.  Since
the representation is unitary, its orthogonal complement
\(Z\) is also invariant.  Thus the orthogonal projection
\(P_Z\) commutes with every \(\rho(g)\) and consequently
with \(\rho(e)\).  It follows that
\[
 u=P_Z\rho(e)jf\in Z\cap\operatorname{im}\rho(e),
 \qquad
 \Theta u=\Theta\rho(e)jf=\tau(a)\Theta jf=\Theta jf.
 \tag{BTH.32}
\]
The rational equality
\[
 Z\cap\operatorname{im}\rho(e)
       =\operatorname{span}_{\mathbb C}
                     j(Z_K\cap\operatorname{im}e)               \tag{BTH.33}
\]
does not require \(e\) to be self-adjoint.  Indeed the
left side equals \(\rho(e)Z\): idempotence makes \(e\)
the identity on its image.  Apply \(\rho(e)\) to the
rational spanning expression for \(Z\); its rational
coefficients and preservation of \(Z\) give the right
side.  The reverse inclusion is immediate.
Equations (BTH.31)--(BTH.33) yield
\(\langle u,jf\rangle=0\).

Write \(u=\Theta^*w\).  Equation (BTH.32) gives
\(\Theta(u-jf)=0\), so
\[
 0=\langle\Theta^*w,u-jf\rangle
   =\langle u,u\rangle-\langle u,jf\rangle
   =\|u\|^2.                                                    \tag{BTH.34}
\]
Hence \(u=0\), and (BTH.32) finishes the proof.
This is a complete proof of the finite assertion in
\texttt{theta\_\allowbreak eq\_\allowbreak zero\_\allowbreak of\_\allowbreak orth}.  In particular the
source's group-algebra level idempotent has not been
silently replaced by an orthogonal projector.
% END BTH9

% BEGIN BTH10
\paragraph{BTH10: the original orthogonal complement is killed, with its exact level normalization.}
\label{claim:BTH10}
Assume (BTH.29), the original \texttt{Ok} block data, and
the hypotheses of BTH4--BTH7.  Then
\[
  \bigl(hform(lH_i)(E_i z,x)=0\text{ for every }z\bigr)
       \quad\Longrightarrow\quad T_i jV_i(f_i(x))=0.
 \tag{BTH.35}
\]

\emph{Proof.}
The source's \(\mathrm{lvlA}(P_0,\psi_0)\) uses
the character \(\psi_0^{-1}\).  Write \(P=P_0\),
\(\psi=\psi_0^{-1}\), and
\[
 p_\psi=\frac1{|P|}\sum_{p\in P}\psi(p^{-1})p,
            \qquad e=Cp_\psi,\qquad a=C.                       \tag{BTH.36}
\]
The central idempotent \(C\) in (BTH.23) commutes with
\(p_\psi\).  The latter is idempotent by reindexing the
double sum, so \(e^2=e\).  If \(f\) is a level-\(\psi\)
form, each summand in \(p_\psi f\) is \(|P|^{-1}f\);
thus \(p_\psi f=f\) and \(ef=af\).

The definition of \(\mathrm{zsOf}\) is the original
rational preimage of the adjoint image.  Consequently
the original good-cut form inclusion has range
\[
       f_i(E_iK^{s_i})
            =Z_{K,i}\cap\operatorname{im}(Cp_\psi).             \tag{BTH.37}
\]
This is a definition-level equality (\texttt{good\allowbreak R\_\allowbreak eq}
and \texttt{range\_\allowbreak to\allowbreak F}), not an equality deduced from
dimensions.  Notice the exact metric factor:
\[
 \mathrm{csum}(f_i(v),f_i(x))
              =|P|\,hform(lH_i)(v,x).                         \tag{BTH.38}
\]
The factor is \(|P|\), because \(lH_i\) was defined with
\(|P|^{-1}\); it is not its reciprocal in (BTH.38).
The model's \texttt{JIso\allowbreak L} identifies the Hilbert pairing
with the \(\sigma\)-image of this rational \(\mathrm{csum}\).
Thus the hypothesis of (BTH.35), together with (BTH.37),
is exactly the orthogonality premise (BTH.31).

BTH6 supplies unitarity, intertwining, and the fact that
\(C\) fixes the theta image.  BTH8 supplies (BTH.30).
BTH9 therefore proves \(\Theta_i(j_i f_i(x))=0\).
Applying \(\mathrm{Tw.out}\) and the actual phase-one
identity gives (BTH.35).
No positive theta-norm formula, full common-image
identity \(\mathrm H5\), or nonbase primitive row is
used in this kernel argument.
% END BTH10

% BEGIN BTH11
\paragraph{BTH11: the stored target basis gives containment and the original rational theta matrix.}
\label{claim:BTH11}
Assume in addition that the supplied bundle's
\texttt{Th\allowbreak Def3} basis is the same \(J_i\) in the reading.
Then the actual good-cut theta image lies in the image of
\(w_{0,i}\), and its original relation matrix is rational
over \(K\).  If
\[
             M_i=(A_3.\mathrm{indefDown}(s0Of(R,F))).\theta(i,t_0),
 \tag{BTH.39}
\]
then, for every \(z\in K^{s_i}\),
\[
                T_i jV_i(f_i(E_i z))=q_iJ_i(M_i z).
 \tag{BTH.40}
\]
There is no additional \texttt{PLev\allowbreak W}, \texttt{LUnif},
or full \(\mathrm H5\) hypothesis.

\emph{Proof.}
The \texttt{Th\allowbreak Def3.li/span} fields say that the complex
vectors \(J_i(e_k)\) are a basis of the original
\(D.Mj(i+1)\).  The map \(q_i\) is injective:
its \texttt{Mj\allowbreak Tie.e\allowbreak W\_\allowbreak norm} formula equates its norm
to a strictly positive fixed multiple of the original
norm, since \(\mathrm{nP}(\mathrm{false})>0\).
The exact source reading is
\[
       w_{0,i}(\sigma(z))
                =\mathrm{Tw.mk}_\epsilon(q_iJ_i z).             \tag{BTH.41}
\]
In particular \(w_{0,i}\) is injective.

Apply BTH7 to \(E_i e_\alpha\).
Its left side belongs to \(q_i(D.Mj(i+1))\), which is
spanned over \(\mathbb C\) by the \(q_iJ_i(e_k)\).
Equation (BTH.41) puts
\(\Theta_i j_i f_i(E_i e_\alpha)\) in
\(\operatorname{im}w_{0,i}\).  Linearity gives this
containment for the entire complexified good-cut
space.  This argument uses the already supplied
\texttt{Th\allowbreak Def3} basis; it does not assume the missing
theta square.

For each column, solve the unique coordinate equation
\[
       w_{0,i}c_\alpha=\Theta_i j_i f_i(E_i e_\alpha).
 \tag{BTH.42}
\]
Applying the tests in (BTH.29) gives \(Ac_\alpha=b_\alpha\),
where every entry of \(A\) and \(b_\alpha\) lies in
\(\sigma(K)\).  The test separation and injectivity of
\(w_{0,i}\) imply that \(A\) has full column rank.
Select a nonzero maximal minor.  Cramer's rule on that
minor gives \(c_\alpha\in\sigma(K)^{r_i}\), including
the empty-dimensional case in the usual unique way.

By definition the source's
\(\mathrm{liftMatC}=\mathrm{relMat}(w_{0,i},tC_i)\).
Its relation is uniquely solvable by (BTH.42), so this
is that matrix of columns \(c_\alpha\).  The operation
\(\mathrm{kMat}_\sigma\) takes their unique preimages
in \(K\).  The matrix in (BTH.39) is exactly that
stored \(\mathrm{liftMat}\), and
\(\mathrm{cr.toL}(z)=f_i(E_i z)\) by the original
coordinate definitions.  Applying (BTH.41) to the
rational relation equation and then undoing the twist
proves (BTH.40).

The source proves a stronger target-space equality
elsewhere using \texttt{PLev\allowbreak W}; the containment above
is sufficient for this base matrix and follows
already from BTH7 and the common basis.
% END BTH11

% BEGIN BTH12
\paragraph{BTH12: the base theta square and its exact remaining producer.}
\label{claim:BTH12}
Under the preceding hypotheses, the original base square is
\[
 \begin{split}
 &D.\theta P\bigl(\mu_0,i+1,
            D.jMt(\mu_0,i+1)(Y.eMt(i,t_0)x)\bigr)\\
 &\hspace{30mm}=J_i(M_i\Pi_i x)
       \qquad(x\in K^{d_i},\ 0\le i<m).
 \end{split}                                                    \tag{BTH.43}
\]
Together with \(\mathrm{RowsHold0}_3\) on this same
bundle, it gives the small-level rank conclusion of SLR8.

\emph{Proof.}
Put \(r=x-E_i\Pi_i x\).
By (BTH.4), \(r\) is \(lH_i\)-orthogonal to
\(\operatorname{im}E_i\).  BTH10 gives
\[
       T_i jV_i(f_i(x))
             =T_i jV_i(f_i(E_i\Pi_i x)).
\]
BTH11 identifies the right side with
\(q_iJ_i(M_i\Pi_i x)\).  BTH7 identifies the left
side with \(q_i\) of the left side of (BTH.43).
The injectivity of \(q_i\) proves (BTH.43).
The identification of \(A_3,\mathrm{zsOf},s0Of\) and of
the maps \(E_i,\Pi_i,J_i,M_i\) with the bundle's stored fields is
required, and follows by substitution once the
same original \texttt{Tw\allowbreak In}/bundle reading is supplied.

Here is a sufficient remaining producer, with the
quantifiers exposed.  For every family
\(\mathcal F\in\mathcal A\), choose one eventual
threshold \(a_\mathcal F\).  For all
\[
  x\ge a_\mathcal F,\qquad x>\mathcal F.x_1,\qquad
           \alpha\in\mathcal F.I_x,                             \tag{BTH.44}
\]
use the \emph{same} original bundle
\(b=\sigma(\mathcal F,x,\alpha)\), where \(\sigma\)
is the common compatible selector.  On it supply:

\begin{enumerate}
 \item the actual \(R,F,Y,Z\) reading, its weak
       characterization and phase-one pins, and its
       identification with \(b.P.A_3,b.P.Zs,b.P.S_0\)
       and the stored \texttt{Th\allowbreak Def3} maps;
 \item for each \(i<m\), a smooth cuspidal
       multiplicity extension of (BTH.10) for the
       finitely many vectors of one basis of the
       original \(H_i\);
 \item for each \(i<m\), one finite separating
       test family (BTH.29), simultaneously rational
       on the full rational source and on the
       original target-coordinate basis;
 \item \(\mathrm{RowsHold0}_3\) on this same bundle.
\end{enumerate}

Items 1--3 give (BTH.43); item 4 and SLR8 give
the base-rank equality.  The elementary passage
from the family threshold to the original
\(\mathrm{SmallLevelRankEq}\) uses the threshold
\(\max(a_\mathcal F,\mathcal F.x_1+1)\).
If this producer is supplied only in the source's
weaker \texttt{Closed} form, its own \(\varepsilon\)
and \(K\) bounds must be retained.  Nothing here
converts such a weaker producer into a universal
one with an arbitrary new bound.
% END BTH12

\subsection*{Removing the finite-action normalization ambiguity}

The next three arguments supply the actual-action assertion used
in BTH5.  The canonical action of CW23 alone would not prove a
statement about a different, already stored
\texttt{Is\allowbreak Weil\allowbreak Datum\allowbreak O}.  The comparison below is on the original
theta functions, with the same additive and splitting characters.

% BEGIN BTH13
\paragraph{BTH13: a continuous adelic unitary character is determined by rational and split points.}
\label{claim:BTH13}
Let \(E/F\) be the original CM extension, \(V\) any
nondegenerate Hermitian \(E\)-space, and
\(\xi:U(V)(\mathbb A_F)\to S^1\) a continuous character.
Suppose that \(\xi\) is trivial on \(U(V)(F)\) and on every
finite element supported at a split base-field place.  Then
\(\xi=1\).  No definiteness hypothesis on \(V\) is required.

\emph{Proof.}
Fix a full lattice \(L\).  For a nonisotropic rational vector
\(v\), its Hermitian projector \(P_v\) is rational, idempotent,
and self-adjoint.  The quasi-reflection
\[
 R_{v,u}=1+(u-1)P_v,\qquad u\bar u=1,
 \quad R_{v,u}^{-1}=R_{v,u^{-1}}                             \tag{BTH.45}
\]
is unitary.  Such reflections generate \(U(V)(F)\).
Here is the argument, including the indefinite case.
If \(x,z\) have the same nonzero norm and \(a=z-x\) is
nonisotropic, put \(A=h(z,a)\), using the convention linear
in the first variable.  Then \(h(a,a)=A+\bar A\ne0\), and
\(u=-\bar A/A\) gives \(R_{a,u}z=x\).
If \(z-x\) is isotropic, \(z+x\) has norm \(4h(x,x)\ne0\);
the preceding construction sends \(z\) to \(-x\), and
\(R_{x,-1}\) sends \(-x\) to \(x\).
A nondegenerate Hermitian space has an orthogonal basis of
nonisotropic vectors: otherwise polarization would make a
nonzero remaining form identically zero.  Fix its first
vector by at most two reflections and induct on its
orthogonal complement.  This is WA10's algebraic
generation proof, now over the rational field \(E/F\).

We first prove triviality on the infinite factor.
Write a rational norm-one \(u\) as
\(u=\alpha/\bar\alpha\) with \(0\ne\alpha\in\mathcal O_E\).
For completeness, choose \(b\) with \(b+u\bar b\ne0\)
and take \(\alpha=b+u\bar b\), then clear its rational
denominator.  Such \(b\) exists because otherwise \(b=1\)
would give \(u=-1\) and then conjugation would fix every \(b\).
For all sufficiently large integers \(n\), put
\[
 N_n=n!,\qquad
 u_n=\frac{1+N_n\alpha}{1+N_n\bar\alpha}.                     \tag{BTH.46}
\]
The numerator and denominator are nonzero, and \(u_n\bar u_n=1\).
At each infinite place \(u_n\to u\); at each fixed finite
place \(u_n\to1\).  At every nonsplit finite place it is a
unit, since conjugation preserves valuation and the two
valuations in (BTH.46) agree.

The projector \(P_v\) has lattice-integral entries outside
finitely many places.  At its finitely many nonsplit
denominator places, \(u_n-1\) and \(u_n^{-1}-1\) tend to zero.
Thus, for all sufficiently large \(n\), the finite component
of \(R_n=R_{v,u_n}\), and its inverse, preserve \(L\) at
every nonsplit place.  Remove its factors at precisely the
finitely many split places where lattice integrality or
inverse integrality fails.  Let \(s_n\) be their product and
let \(k_n\) be the remaining finite component.  Then
\[
 R_n=R_{n,\infty}\,s_n\,k_n,\qquad
 k_n\in K_L:=\prod_{w<\infty}U(L_w).                         \tag{BTH.47}
\]
Each \(U(L_w)\) is a closed subgroup of
\(\operatorname{GL}(L_w)\), hence compact.  The product
\(K_L\) has its product topology as a compact open
subgroup of the finite adelic group.
For any fixed finite place \(w\), \(R_n\to1\) there;
eventually its component is integral with integral inverse,
so no factor is removed there and \((k_n)_w\to1\).
Consequently \(k_n\to1\) inside this one fixed compact
product.  This is the required restricted-product
convergence, not merely componentwise convergence of
unbounded adelic elements.

Rational triviality gives \(\xi(R_n)=1\), and split
triviality gives \(\xi(s_n)=1\).  Continuity in (BTH.47)
therefore proves
\(\xi(R_{v,u,\infty})=1\).
Reflection generation proves the same for the infinite
component of every rational unitary matrix.

Those components are dense in the full infinite unitary
group.  More generally, rational unitary matrices satisfy
weak approximation at any finite set of places.  To see
this, multiply all the desired local matrices by one
rational norm-one scalar \(z\) so that \(1+zg_w\) is
invertible at each selected place.  There are only
finitely many forbidden scalars; the rational norm-one
scalars are infinite, for instance by the ratios
\((t+\beta)/(t+\bar\beta)\) with \(t\in F\) and
\(\beta\notin F\).
The Cayley transform
\[
 X=(1-zg)(1+zg)^{-1},\qquad
 X^*=-X,\qquad zg=(1-X)(1+X)^{-1}                            \tag{BTH.48}
\]
identifies this open set with an open subset of the
\(F\)-vector space of skew-Hermitian matrices.
Field weak approximation applied to its coordinates,
and openness of the nonzero determinant conditions,
give a rational approximation; undo \(z\).
This proof applies to indefinite forms as well.
Continuity now gives \(\xi=1\) on the entire infinite
factor.

For the finite factor we prove the required approximation
with a single fixed lattice.  Given a finite adelic
unitary \(g\) and an integer \(N>0\), choose a finite set
\(S\) containing all primes over \(N\) and every place
where \(g\) or \(g^{-1}\) fails to preserve \(L\).
By (BTH.48), choose rational \(\gamma_0\) so that
\(\gamma_0g\) is integral and invertible at \(S\) and is
congruent to \(1\) modulo \(NL_w\) at the primes over \(N\).
Factor \(\gamma_0\) into the rational reflections (BTH.45).

Correct each reflection parameter while preserving this
precision at \(S\).  The elementary ratio
\(\alpha/\bar\alpha\), together with field weak
approximation on nonzero local choices of \(\alpha\),
approximates any prescribed norm-one parameters at
finitely many places.  Choose the corrected parameter
close to its old value at \(S\), and close to \(1\) at
every other nonsplit denominator place of its projector.
Every rational norm-one parameter has valuation zero at
a nonsplit prime.  Hence the corrected reflection and
its inverse preserve \(L\) at every nonsplit place
outside \(S\); the closeness to \(1\) clears the finitely
many projector denominators.  Choose enough precision
at \(S\) that the corrected product \(\gamma\) retains
the original integrality and congruence there.
This is QD13's correction argument; all of its ingredients
have just been checked without a definiteness assumption.

The element \(\gamma g\) can now fail lattice
stabilization only at finitely many split places outside
\(S\).  Let \(s\) equal its components there and \(1\)
elsewhere; let \(k\) equal \(1\) there and its components
elsewhere.  These support cuts are unitary, commute, and
give the exact factorization
\[
 g=\gamma_f^{-1}s k,\qquad k\in K_L,\qquad
       (k-1)L_w\subset NL_w\quad(w\mid N).                   \tag{BTH.49}
\]
This proves the needed generalization of QD14.
Take \(N=n!\).  The resulting \(k_n\) converge to \(1\)
in \(K_L\).
The rational identity and the already proved infinite
triviality give \(\xi(\gamma_f)=1\);
split triviality gives \(\xi(s)=1\).
Thus \(\xi(g)=\xi(k_n)\to1\), proving finite triviality.
The infinite and finite factors together prove the claim.
These are the original quantifiers of
\texttt{unitary\_\allowbreak adelic\_\allowbreak character\_\allowbreak trivial\_\allowbreak rat\_\allowbreak and\_\allowbreak split}.
% END BTH13

% BEGIN BTH14
\paragraph{BTH14: full Weil data with the same characters agree on the original theta space.}
\label{claim:BTH14}
Let \(\omega,\omega'\) be two full
\texttt{Is\allowbreak Weil\allowbreak Datum\allowbreak O} actions for the same
\(D,\psi,S,T,\chi_S,\chi_T\), on the original functions.
Then for every theta function \(f\) and every adelic
pair \(p\),
\[
                         \omega(p)f=\omega'(p)f.             \tag{BTH.50}
\]
In particular this applies to the self-dual lattices
and the already stored action used in BTH5.

\emph{Proof.}
The operator \(\omega'(p^{-1})\omega(p)\) preserves
theta functions, commutes with every Heisenberg
translation by the two covariance laws, and preserves
the original Haar norm by the two \texttt{om\_\allowbreak unit}
laws.  CW15's scalar lemma, including its norm
hypothesis, therefore gives
\[
          \omega(p)f=c(p)\omega'(p)f,\qquad |c(p)|=1.          \tag{BTH.51}
\]
If the theta space is zero there is nothing to prove.
Otherwise fix a nonzero theta function \(f_0\).
The operator \(\omega'(p)\) is invertible on that
space, so \(c(p)\) is unique.  Applying the two group
laws to \(f_0\) proves that \(c\) is a character.
It is continuous: near any \(p_0\), choose \(w\) with
\((\omega'(p_0)f_0)(w)\ne0\).  The quotient
\[
       c(p)=\frac{(\omega(p)f_0)(w)}
                        {(\omega'(p)f_0)(w)}
\]
there is continuous by the literal pointwise
\texttt{om\_\allowbreak cont} laws.
Translate \(f_0\) by a Heisenberg operator if needed
so that \(f_0(0)\ne0\).  Both rational
\texttt{om\_\allowbreak aut} laws evaluated at \(0\) then give
\(c(\gamma,\gamma')=1\) for every rational pair.

For a split-supported pair \(p\), choose the same
determinant lifts \(x,y\) in both split clauses and
write \(z=\chi_T(x)^{-1}\chi_S(y)\).
They supply positive real numbers \(r,r'\) with
\[
 \ell(\omega(p)f)=rz\ell(f),\qquad
 \ell(\omega'(p)f)=r'z\ell(f)                                \tag{BTH.52}
\]
for every functional invariant under the specified
full split half.  There are such \(\ell,f\) with
\(\ell(f)\ne0\).
Here is why the earlier coinvariant proof applies.
The finite character projectors of CW3, on the
standard self-dual finite Hom lattice, express any
nonzero theta vector as a finite sum of translates
of lattice-fixed vectors, so at least one is nonzero.
The required self-dual lattices are the images of
the standard lattices already stored in \(S,T\)
under their frames; their duals transform by the
same isometry, as checked in CW24.
On each of the two split half-lattices the corrected
Heisenberg action in CW3 equals the raw action:
the phase is zero by isotropy of that half.
WA5's algebraic coinvariant argument applied to
this vector gives the required functional and
theta vector.  This construction concerns only
Heisenberg operators and lattices, not either
choice of Weil action.
Substitute (BTH.51) into (BTH.52) and cancel
the nonzero \(z\ell(f)\): \(c(p)r'=r\).
Since \(|c(p)|=1\) and \(r,r'>0\), it follows
that \(r=r'\) and \(c(p)=1\).

Restrict \(c\) separately to \(U(S)(\mathbb A_F)\)
and \(U(T)(\mathbb A_F)\).  Each restriction meets
all hypotheses of BTH13, so both are trivial.
The decomposition \(p=(p_S,1)(1,p_T)\) gives
\(c(p)=1\), proving (BTH.50).
The assertion is equality on theta functions;
no equality of arbitrary extensions to the ambient
space of all functions has been asserted.
This proves the source's
\texttt{weil\allowbreak Datum\_\allowbreak unique\_\allowbreak on\_\allowbreak theta\allowbreak Functions}.
% END BTH14

% BEGIN BTH15
\paragraph{BTH15: the original action has the exact finite lattice-stabilizer character.}
\label{claim:BTH15}
Use the original self-dual lattices \(L_S,L_T\),
their frames, a supplied \(\eta:\texttt{Is\allowbreak Eta\allowbreak Tilde}\),
and any full original
\(\omega:\texttt{Is\allowbreak Weil\allowbreak Datum\allowbreak O}(\psi,S,T,\chi_S,\chi_T)\).
Suppose \(\chi_S,\chi_T\) are unramified away from
the finitely many split pairs selected by an idempotent
\(e\).  If \(g,h\) are finite lattice stabilizers and
\(f\) is an original theta function satisfying the
original \texttt{Is\allowbreak Lat\allowbreak Vac} law for the Hom lattice
\(\Lambda\), then
\[
 \omega(g,h)f=
 \chi_T(\det\nolimits_e g)^{-1}
          \chi_S(\det\nolimits_e h)\,f.                       \tag{BTH.53}
\]
Here \(\det_e g=e\det g+1-e\), with the source's
idele interpretation; the two characters are crossed.

\emph{Proof.}
First let the pair be supported at one split place.
Its action preserves \(\Lambda\) in both directions
and commutes with its split-half projections.
Choose the literal determinant lifts \(x,y\), and
let \(r>0\) be the number supplied by
\(\texttt{om\_\allowbreak split}\), with
\(z=\chi_T(x)^{-1}\chi_S(y)\).
The theta function
\[
                         d=\omega(g,h)f-rzf
\]
is fixed by both complementary split half-lattices:
this holds for \(f\) by \texttt{Is\allowbreak Lat\allowbreak Vac}, and for
\(\omega(g,h)f\) by covariance and lattice
stabilization.  If \(d\ne0\), the WA5 coinvariant
proof applied to \(d\) produces a linear functional
\(\ell\), invariant under the full required half,
with \(\ell(d)\ne0\).
Self-duality supplies the separation used in that
proof; finite-level smoothness supplies the
telescoping sums.  The same
\(\texttt{om\_\allowbreak split}\) clause applied to \(f\)
gives \(\ell(d)=0\), a contradiction.
Thus \(\omega(g,h)f=rzf\).

If \(f=0\) the desired identity is immediate.
Otherwise its continuous norm-square descends to
the compact rational additive quotient of the
Hom adelic space.  Its Haar integral is finite
and positive: compactness bounds it, and a point
where it is nonzero has an open neighborhood of
positive Haar measure.  The norm identity
\(\texttt{om\_\allowbreak unit}\) now gives
\(r^2\|f\|_2^2=\|f\|_2^2\); since \(r>0\),
we obtain \(r=1\).  This proves the exact split
normalization for the supplied action itself.

Next let the pair be supported at a nonsplit place
\(v\), with both characters unramified there.
CW23 constructs a canonical full action \(\omega_0\)
on these same spaces and lattices, with characters
\(\eta^{\dim S},\eta^{\dim T}\) and its proved
local fixed-vacuum law.  Put
\[
 \nu_S=\eta^{-\dim S}\chi_S,\qquad
 \nu_T=\eta^{-\dim T}\chi_T.
\]
They are Hecke characters trivial on the ideles of
\(F\), by the exact rank restrictions.
The determinant-root construction WA18 gives
\[
 \omega_1(g,h)=\widetilde\nu_T(g)^{-1}
                  \widetilde\nu_S(h)\omega_0(g,h),             \tag{BTH.54}
\]
a full datum with the same characters as the
given \(\omega\).
BTH14 therefore identifies \(\omega\) and
\(\omega_1\) on every original theta function.

The scalars in (BTH.54) are one at this pair.
Indeed a norm-one determinant \(a\) at \(v\) is
a unit, since conjugation preserves valuation.
The original host's \texttt{cm\allowbreak Unram} condition
makes the involution nontrivial on the residue
field at every conjugation-fixed prime.
Choose an integral \(b\) for which
\[
                  x=b+a\bar b
\]
has nonzero residue.  Such \(b\) exists: if the
expression vanished for every residue \(b\), then
\(b=1\) would give \(a=-1\) and the remaining
identities would make the residue involution
trivial.  Thus \(x\) is a unit and
\(x=a\bar x\), so \(a=x/\bar x\).
Take this determinant root at \(v\) and \(1\)
elsewhere.  Both \(x\) and \(x^{-1}\) are
integral finite ideles, supported at \(v\).
The original characters and \(\eta\) are
unramified there, so their values on \(x\)
are one.  The same construction applies to
the other determinant, proving the assertion.

The canonical local law asks for fixing under
every Heisenberg vector in the local lattice.
The original \texttt{Is\allowbreak Lat\allowbreak Vac} first gives fixing
under \(2u\).  The host makes every dyadic
place split, so \(2\) is a unit at the present
nonsplit \(v\).  Divide any local lattice vector
by \(2\), still inside that lattice, to obtain
the required stronger local fixing.  Thus the
canonical \texttt{Fixes\allowbreak Vac\allowbreak At} applies, and
(BTH.54) gives \(\omega(g,h)f=f\) for this
original action.

For the general finite stabilizer pair, first
truncate to finitely many places.
At a selected split pair the first part gives
exactly its determinant factor in (BTH.53).
At an unselected split pair both determinant
lifts are integral units: a matrix and its
inverse preserve the framed lattice, so its
determinant and inverse determinant are integral.
The two unramified characters therefore give
factor one.  At each nonsplit place the preceding
argument gives factor one.  Multiplication proves
(BTH.53) for every finite-support truncation.
All these truncations lie in the same compact
product of framed lattice stabilizers and
converge to the given pair.  The original
\texttt{om\_\allowbreak cont}, the continuity of the
characters, and the continuous determinant cuts
let us pass to the limit at every function
argument.  This proves (BTH.53).

The precise source chain is recorded in the following files:
\begin{quote}\small
\texttt{WeilDatumUniqueOnThetaFunctions};\\
\texttt{WeilDatumFixesVacuumNonSplitPlaceUnramified};\\
\texttt{WeilOperatorLatticeStabiliserFixesVacuumNonsplit};\\
\texttt{SplitPlaceWeilRepLatticeVacuumCharacter};\\
\texttt{WeilLatticeVacuumStabilizerSplittingCharacters}
(\texttt{stab\_vacuum}).
\end{quote}
The canonical-to-original comparison is now proved;
it is not an additional finite-covariance producer.
% END BTH15

\subsection*{What this closes, and what it leaves open}

The mathematical arguments BTH1--BTH15 are category B
proofs of their stated implications.  In particular,
BTH3 closes the actual integral-convergence guard
once the original models and weak Weil data are
supplied.  BTH4 constructs the good-projected Hilbert
lift, rather than leaving a second lift-existence
problem.  BTH6 reconstructs the required target
unitarity and base equivariance.  BTH10 is the full
finite orthogonal-kernel proof, and BTH11 replaces
the unnecessarily strong common-image package by
a consequence of the same stored basis.
BTH13--BTH15 remove the scalar ambiguity between
the canonical Weil construction and the original
action, proving its exact finite stabilizer character.

The remaining existence obligations are not classified
as B merely because the resulting formulas are short.
The precise residual table is:

\begin{center}
\begin{tabular}{p{0.18\linewidth}p{0.60\linewidth}p{0.10\linewidth}}
\hline
Obligation & Required original data & Class\\ \hline
Base reading &
Actual \texttt{Read}/\texttt{Read\allowbreak L}, their
\texttt{Lift\allowbreak Tie.theta}, \texttt{mt\allowbreak Read}, \texttt{mj\allowbreak Read},
weak models and pins, identified with the selected
bundle's row family and basis.  An unrelated model
with equal dimensions is insufficient. & C\\
Base extension &
For a basis of each original \(HV(i,t_0)\), extend
its now-convergent lowest-type integral to
\(\mathrm{Wall}_i\).  This includes the actual
smooth, equivariant, cuspidal multiplicity property. & C\\
Base rational tests &
The simultaneous finite separating tests (BTH.29)
for the exact field, embedding, theta map, and
target basis. & C\\
Base rows &
\(\mathrm{RowsHold0}_3\) on the same bundle;
its primitive row construction is a separate
producer. & C\\ \hline
\end{tabular}
\end{center}

The first three entries refine the theta part of the
previous joint same-bundle producer.  They are
shared mathematical existence tasks, not three
independent freely choosable witnesses: every
entry uses the same \(R,F,Y,Z,\omega_i,io_i,pt_i\),
measures, field, and stored weak choices.  There
is no new category-D instance ambiguity in the
displayed implication once these identifications
are supplied.  Their actual construction remains
unclosed, and the original joint producer must
not be marked complete.

For the remaining extension, the pinned source gives a
concrete next mathematical target.  In
\texttt{Exists\allowbreak Theta\allowbreak Lift\allowbreak Linear\allowbreak Map\allowbreak Cusp\allowbreak Form\allowbreak Multiplicity},
the declaration \texttt{firesAt}
proves the pointwise extension by first establishing
smoothness, uniform moderate growth, zero transvection
constant terms, compact-type equivariance, and annihilation
by the lowering operators of the holomorphic discrete
series.  Its \texttt{exists\_\allowbreak extension} then realizes
these functions in the prescribed multiplicity model.
The exact suppliers include:
\begin{quote}\small
\texttt{thetaLift\_definite\_smooth\_moderate\_growth\_annihilated};\\
\texttt{thetaLift\_definite\_constantTerm\_eq\_zero};\\
\texttt{fock\_vectors\_killed\_heisenberg\_average\_isotropic};\\
\texttt{boundedHolomorphicForm\_eq\_multiplicitySpace\_value}.
\end{quote}
Their substantial analytic and representation-theoretic
proofs have not been replaced by naming a published
theorem.  The source's use of
\texttt{Published.Cusp\allowbreak Rapid\allowbreak Decay} in that route is
likewise not a new category-A certificate in this chapter.

The reduced base route does not need the remaining
fields of \texttt{Ex\allowbreak At}, the chain lift \(\mathrm{Th}\),
an independent \(\mathrm{thG}\) existence theorem,
a separate \texttt{Lift\allowbreak Test}, nonbase test families,
\texttt{PLev\allowbreak W}, \texttt{LUnif}, \(\mathrm H5\),
a positive theta-norm formula, \texttt{Rows\allowbreak Hold\allowbreak Lam3},
or the all-tuple \texttt{Tup\allowbreak All\allowbreak Used} comparison.
This is a reduction of the sufficient producer for
this consumer.  It does not establish existence
of the shared selector, the original analytic
reading, or the full RH/GRH proof.
